% Замените поясняющий текст собственными результатами.
% В шаблоне нет готового решения и обязательных внешних изображений.
\reportelement{Введение}

Сформулируйте цель и задачи лабораторной работы. Перечислите используемые
инструменты; персональные параметры укажите в таблице~\ref{tab:parameters},
изменив набор строк под своё задание.

\begin{table}[H]
  \caption{Персональные параметры}\label{tab:parameters}
  \centering
  \begin{tabularx}{\textwidth}{|p{0.46\textwidth}|X|}
    \hline
    Параметр & Значение\\ \hline
    Номер ИСУ & \StudentISU\\ \hline
    Персональный порт & 5489\\ \hline
    Персональный префикс & matveev\\ \hline
    Окружение & linux/x86\_64\\ \hline
    Секретный ключ & ------\\ \hline
  \end{tabularx}
\end{table}

Результат хранится в ветке \texttt{lab-\LabNumber} своего форка:
\url{\RepositoryURL}.

\chapter{Ход выполнения}
\section{Первый рабочий образ}

На рисунке~\ref{fig:code1} сделал Fork репозитория курса на GitHub и клонировал свой форк. Положил приложение в solutions/assist-py (без истории git оригинала).

\reportscreenshot{image.png}{Результат этапа}{fig:code1}

Создал файл Dockerfile в корне проекта на рисунке~\ref{fig:code2}.

\reportscreenshot{Screenshot_1.png}{Dockerfile}{fig:code2}

Собрал образ на рисунке~\ref{fig:code3}.
\reportscreenshot{Screenshot_2.png}{Сборка}{fig:code3}

Запустил контейнер с переменными окружения на рисунке~\ref{fig:code4}.
\reportscreenshot{Screenshot_3.png}{Запуск}{fig:code4}

Проверил работоспособность на рисунках ~\ref{fig:code5}-~\ref{fig:code9}.
\reportscreenshot{Screenshot_4.png}{Проверка}{fig:code5}
\reportscreenshot{Screenshot_5.png}{Проверка}{fig:code6}
\reportscreenshot{Screenshot_6.png}{Проверка}{fig:code7}
\reportscreenshot{Screenshot_7.png}{Проверка}{fig:code8}
\reportscreenshot{Screenshot_8.png}{Проверка}{fig:code9}

\section{Оптимизация образа}

Создан новый файл Dockerfile.optimized изображенный на рисунке ~\ref{fig:code10}.

В отличие от Dockerfile из этапа 1, где образ собирался в один этап, новая версия использует multi-stage сборку: зависимости устанавливаются в промежуточном образе builder, а в финальный образ runtime переносятся только готовые пакеты и код приложения. За счёт этого в итоговый образ не попадают артефакты установки.

Установка зависимостей выполняется командой pip install --no-cache-dir --prefix=/install -r requirements.txt. Флаг --no-cache-dir отключает сохранение кэша pip, поэтому скачанные пакеты не увеличивают размер слоя. Параметр --prefix=/install позволяет установить пакеты в отдельный каталог, который затем целиком копируется в финальный образ командой COPY --from=builder /install /usr/local.

\reportscreenshot{Screenshot_9.png}{Dockerfile.optimized}{fig:code10}

Размеры образов изображены на рисунке ~\ref{fig:code11}. Размер образа v1 составил 188 МБ, образа v2 составил 164 МБ, то есть образ уменьшился на 24 МБ. Основной вес образа определяется базовым образом python:3.11-slim и самими зависимостями, поэтому разница небольшая, multi-stage сборка убирает только артефакты установки, но не сами пакеты.

\reportscreenshot{Screenshot_10.png}{Сравнение}{fig:code11}

Ответ на вопрос о кэшировании. Копирование requirements.txt отдельно от кода ускоряет повторные сборки, потому что Docker кэширует каждый слой и сбрасывает кэш только тогда, когда изменились инструкция или копируемые ею файлы. При изменении бизнес-логики файл requirements.txt остаётся прежним, поэтому слои с копированием списка зависимостей и установкой пакетов берутся из кэша, а пересобирается только слой с кодом приложения. Если бы весь проект копировался до установки зависимостей, любое изменение кода сбрасывало бы кэш, и пакеты скачивались бы заново при каждой сборке.

\section{.dockerignore - чистота в образе}

В корне проекта создан файл .dockerignore с содержимым изображенным на рисунке ~\ref{fig:code12}. Также был пересобран образ.

\reportscreenshot{Screenshot_11.png}{.dockerignore файл}{fig:code12}

Без .dockerignore в образ могли попасть все файлы из контекста сборки, которые затрагивает инструкция COPY. В этом случае это:
\begin{enumerate}
    \item каталоги \texttt{\_\_pycache\_\_/} и файлы *.pyc, *.pyo (скомпилированный байт-код, который бесполезен в образе и создаётся под конкретную среду разработки);
    \item каталог .git со всей историей репозитория;
    \item .gitignore, README.md, Dockerfile* (служебные файлы, которые приложению не нужны);
    \item виртуальные окружения .venv/ и venv/ (могут весить сотни МБ и содержать пакеты, собранные под другую систему);
    \item файл .env с паролями, токенами и секретными ключами;
    \item лог-файлы *.log, в которых могут быть адреса, данные пользователей и внутренняя информация.
\end{enumerate}

Всё это увеличивает размер образа и замедляет сборку, потому что весь контекст передаётся демону Docker, а любое изменение таких файлов сбрасывает кэш слоя с COPY.

Наличие .git в образе является потенциальной уязвимостью, потому что в этом каталоге хранится полная история проекта. Любой, кто получил доступ к образу (например, из публичного реестра) или к контейнеру, может выполнить git log и git show, восстановить исходный код и увидеть все прошлые версии файлов. Даже если пароль, ключ или токен был удалён из кода одним из последующих коммитов, он остаётся в истории и его можно достать. Также в .git видны адреса удалённых репозиториев, имена и почты авторов, а иногда и данные для доступа в конфигурации. Помимо риска утечки, .git увеличивает размер образа.

\section{Healthcheck - готовность к оркестрации}

Добавлена проверка здоровья, в конец Dockerfile.optimized (после EXPOSE) добавлен код изображенный на рисунке ~\ref{fig:code13}.

\reportscreenshot{Screenshot_12.png}{Healthcheck}{fig:code13}

Собрал финальную версию образа, проверка статуса здоровья изображена на рисунке ~\ref{fig:code14}.

\reportscreenshot{Screenshot_13.png}{Проверка}{fig:code14}

После добавления установки \texttt{curl} размер образа вырос на 14 МБ (рисунок ~\ref{fig:code15}).
Рост связан с тем, что вместе с \texttt{curl} через \texttt{apt-get} устанавливаются его зависимости (\texttt{libcurl}, библиотеки для TLS, сжатия и др.).

\reportscreenshot{Screenshot_14.png}{Размеры}{fig:code15}

Проверку можно выполнить средствами самого Python, который уже есть в образе:

\begin{verbatim}
HEALTHCHECK --interval=30s --timeout=5s --start-period=10s
--retries=3 \
CMD python -c "import os,urllib.request as u; \
u.urlopen('http://localhost:'+os.environ['PORT']+'/health')" ||
exit 1
\end{verbatim}

Это не требует установки дополнительных пакетов и не увеличивает размер образа.

\section{Финальный артефакт и тегирование}

Добавил семантические теги. Артефакт выполнения изображен на рисунке ~\ref{fig:code16}.

\reportscreenshot{Screenshot_15.png}{Артефакт}{fig:code16}



\reportelement{Заключение}

В ходе работы собран образ приложения на FastAPI и пройден путь от базового \texttt{Dockerfile} до образа, готового к оркестрации. Образ \texttt{matveev-assistant:v1} запущен с персональным портом 5489 и проверен запросами к эндпоинтам. В \texttt{Dockerfile.optimized} применены multi-stage сборка, отключение кэша pip и порядок слоёв, ускоряющий пересборку. Файл \texttt{.dockerignore} исключил из контекста лишние файлы. Добавленный \texttt{HEALTHCHECK} обеспечил статус \texttt{healthy}, а теги версии и даты сделали образ воспроизводимым. Полученные навыки станут основой для работ по Docker Compose и Kubernetes.

% Оставьте только использованные источники и цитируйте их в тексте.
\begin{thebibliography}{9}
  \bibitem{docker-docs}
  Docker Documentation: документация Docker.
  URL: \url{https://docs.docker.com/}
  (дата обращения: 9.11.2026).
\end{thebibliography}
